
\section*{Part IV. Applications of atomic physics \hfill 13 / 15}
\subsection*{1. Optical transitions: identify the allowed operator}
An electromagnetic transition has both angular and parity constraints.
Electric multipole \(Ek\): rank \(k\), parity \((-1)^k\).
Magnetic multipole \(Mk\): rank \(k\), parity \((-1)^{k+1}\).
For a transition \(J,m\to J',m'\), angular momentum requires
\[
\boxed{|J-J'|\le k\le J+J',\qquad
m'=m+q,\ |q|\le k.}
\]
\begin{center}\small
\begin{tabular}{@{}llll@{}}\toprule
Type & Parity & Usual \(\Delta J\) & Single-electron orbital check\\\midrule
E1 & changes & \(0,\pm1\), not \(0\leftrightarrow0\) & \(\Delta l=\pm1\)\\
M1 & same & \(0,\pm1\), not \(0\leftrightarrow0\) & typically \(\Delta l=0\)\\
E2 & same & \(0,\pm1,\pm2\), triangle rule & \(\Delta l=0,\pm2\), triangle rule\\\bottomrule
\end{tabular}
\end{center}
Total \(L,S\) rules and level mixing depend on the coupling
approximation. \emph{Example:} \(s\to p\) can be E1;
\(s\to d\) can be E2; \(s\to p\) cannot be E2
because the orbital parity changes.

\subsection*{2. Einstein coefficients, linewidths, lifetimes}
Stimulated absorption/emission depend on the incident spectral
energy density; spontaneous decay occurs even in vacuum.
For an E1 transition between specified states (with
consistent treatment of polarisation/sublevels),
\[
\boxed{A_{if}=
\frac{\omega_{if}^3}{3\pi\epsilon_0\hbar c^3}
|\langle f|\mathbf d|i\rangle|^2,\quad
\tau_i=\Big(\sum_{f:E_f<E_i}A_{if}\Big)^{-1}.}
\]
Here \(\mathbf d=-e\mathbf r\).
For \(R(r)\) normalisation, the radial dipole integral is
\(\int_0^\infty R_f^*(r)R_i(r)r^3\,dr\);
multiply by the angular factor or reduced matrix element.
The natural linewidth (angular frequency) is of order
\(\Gamma=\sum_fA_{if}=1/\tau\).
With \(\hbar\omega\) in atomic units, convert
energies, dipoles and time consistently
(\(t_{\rm au}=2.418884\times10^{-17}\ {\rm s}\)).

\subsection*{3. Laser cooling and trapping}
Absorbing a photon transfers momentum \(+\hbar\mathbf k\);
spontaneous emission gives a random recoil.
For a driven two-level atom, the scattering rate is
\[
R_{\rm sc}=\frac{\Gamma}{2}
\frac{s}{1+s+(2\Delta/\Gamma)^2},\quad
\mathbf F_{\rm sc}=\hbar\mathbf k R_{\rm sc},\quad
s=I/I_{\rm sat}.
\]
The atom experiences Doppler-shifted detuning
\(\Delta_{\rm eff}=\Delta-\mathbf k\cdot\mathbf v\).
Opposing red-detuned beams yield velocity-dependent damping
(optical molasses). A basic two-level Doppler limit is
\(T_D=\hbar\Gamma/(2k_B)\). In a magneto-optical trap
a magnetic-field gradient plus polarisation-dependent
Zeeman shifts adds a restoring force.
\begin{methodbox}
\textbf{Method:} check selection rules first;
obtain matrix elements and rates, then lifetime.
For cooling problems, identify detuning, Doppler shift,
scattering momentum and restoring/damping force.
\end{methodbox}
\src{Week 9 pp.~8--21; Week 10 pp.~4--58 (radiation), 60--130 (optical pumping/cooling).}
